% Fig. 4 — H4 timeline. Used by sections/06-evaluation.
% Three lanes: the transport view (identical in both cases), an unprepared
% command path, a prepared one. Timings come from the bench comparison
% (n = 3 per arm) recovered from git commit 08bfc34; see review/REPONSE.md.
%
% The two outcomes are distinguished by marker SHAPE as well as colour
% (crossed box for the aborted path, filled square for a state transition),
% so the figure still reads when printed in greyscale.
\begin{figure}[!t]
\centering
\resizebox{\columnwidth}{!}{%
\begin{tikzpicture}[x=1cm, y=1cm, font=\scriptsize]
  \def\L{7.8}
  \tikzset{
    lane/.style={draw=figslate, line width=0.4pt,
                 -{Stealth[length=3pt,width=2.4pt]}},
    lanename/.style={font=\scriptsize, align=right, text=black},
    ev/.style={draw, line width=0.4pt, minimum size=0.19cm, inner sep=0pt},
    note/.style={font=\tiny, text=figslate},
  }

  % ---- lane A : what the transport layer reports -------------------------
  \draw[lane] (0,2.15) -- (\L,2.15);
  \node[lanename, left] at (-0.15,2.15) {Transport\\view};
  \node[ev, fill=figblue,     draw=figblue]      (req) at (0.65,2.15) {};
  \node[ev, fill=figblue!30,  draw=figblue]      (ack) at (1.20,2.15) {};
  \node[note, above] at (0.65,2.28) {request};
  \node[note, above right] at (1.32,2.26)
    {acknowledgment: identical whether or not the robot moves};

  % ---- lane B : command issued without the preparation sequence ----------
  \draw[lane] (0,1.05) -- (\L,1.05);
  \node[lanename, left] at (-0.15,1.05) {Unprepared\\path};
  \draw[dashed, draw=figslate!60, line width=0.4pt] (1.20,0.93) -- (1.20,1.17);
  \draw[decorate, decoration={brace, amplitude=3pt}, draw=figslate!70]
        (1.20,1.20) -- (6.35,1.20);
  \node[note] at (3.8,0.80) {no state transition, pose unchanged};
  \node[ev, draw=figrust, fill=figrust!15, minimum size=0.24cm] (ab) at (6.75,1.05) {};
  \draw[figrust, line width=0.5pt] (ab.south west) -- (ab.north east)
                                   (ab.north west) -- (ab.south east);
  \node[note, right] at (6.95,1.05) {abort, \SI{150}{\second}};

  % ---- lane C : command issued after the preparation sequence ------------
  \draw[lane] (0,-0.05) -- (\L,-0.05);
  \node[lanename, left] at (-0.15,-0.05) {Prepared\\path};
  \node[ev, fill=figteal, draw=figteal] at (1.35,-0.05) {};
  \node[ev, fill=figteal, draw=figteal] at (1.95,-0.05) {};
  \node[ev, fill=figteal, draw=figteal] at (5.70,-0.05) {};
  \node[note, above] at (1.35,0.08) {ready};
  \node[note, above] at (2.05,0.08) {navigating};
  \node[note, above] at (5.70,0.08) {arrived};
  \draw[decorate, decoration={brace, mirror, amplitude=3pt}, draw=figslate!70]
        (1.95,-0.20) -- (5.70,-0.20);
  \node[note] at (3.8,-0.48) {pose advances continuously};
  \node[note, right] at (5.90,-0.48) {median \SI{40.9}{\second}};

\end{tikzpicture}%
}
\caption{The transport lane is identical whether or not the command executes. Only an
independent state topic distinguishes the two lower cases. Timings from the bench comparison of
Section~\ref{sec:eval}-C ($n=3$ per arm).}
\label{fig:h4}
\end{figure}
